\begin{table}[htbp]
\centering
\footnotesize
\caption{\textbf{6.4} Los campeones: clientes con puntaje 4 o 5 en las tres dimensiones.}
\label{tab:p6_campeones}
\begin{tabular}{lr}
\toprule
\textbf{concepto} & \textbf{valor} \\
\midrule
clientes campe\'ones (puntaje 4 o 5 en R, F y M) & 783 \\
fracci\'on de la base de clientes & 21.41\,\% \\
venta que generaron (\pounds) & 1\,302\,057.01 \\
venta total de la ventana, todos los clientes con ID (\pounds) & 2\,432\,839.69 \\
\textbf{fracci\'on de la venta de la ventana} & \textbf{53.52\,\%} \\
\bottomrule
\end{tabular}
\\[0.3em]{\scriptsize \textbf{Qu\'e leer:} el 21.4\,\% de los clientes genera el 53.5\,\% de la venta del trimestre. Es casi exactamente la ley de Pareto, y es el argumento comercial entero del RFM en una l\'inea.}
\end{table}
